\section{The split-prime Hecke identity for the actual twisted tests}

% BEGIN P75H1
\paragraph{P75H1. A common value for the twists and their Hermitian shifts.}
Write $\mathcal O=\mathcal O_E$. Fix an admitted family $\phi$ at
$(M,R)$, a degree-$N$ column lattice $X$ agreeing with
$X_0=\mathcal O^N$ at $RM$, and a good prime $P$ of norm $q$
for $RM$. Write $Q=\bar P$ and choose $b\in Q$ with $b\equiv1\pmod P$.
In the following formulas $\Lambda$ is a definite rank-$n$ lattice of
the given level. At a fixed matrix $x$ put $u(v)=Rxv$.
Enlarge $X$ by a line in $X/PX$ at $P$ to give the lattices $X_i^+$,
and take the inverse images $X_i^-$ of the hyperplanes in $X/QX$.
The lists have the same length $A=(q^N-1)/(q-1)$ and need no
specified pairing. Pick $\eta_i\in X\setminus X_i^-$ and representatives
$\xi_j$ of $X/PX$. Use the rational Hermitian shifts
\[
 \beta_{ij}=a_i d_j^*+d_j a_i^*,\qquad
 a_i=R\eta_i,\quad d_j=Rb\xi_j .
\]

There is one positive integer $s\equiv1\pmod{R^2}$ with
$sX_0\subset qX$. In fact an integer prime to $RM$ clears $X_0$ into
$X$, and $q$ is prime to $R$. Take a multiple of their product
which is one modulo $R^2$. Thus $sX_0$ lies in every $X_i^\pm$.
All twists at $\Lambda$ in these lists take the value
$c=\phi_\Lambda(sx)$ on their support and zero elsewhere, with support
condition $u(Y)\subset\Lambda$ for the twisting lattice $Y$
(the value on that set can also be zero).
To verify independence of the chosen multiplier in the definition of
the twist, let $s'$ be another permitted choice for $Y$. Then
$(s-s')X_0/R^2\subset Y$: away from $R$ division by $R^2$ is a unit,
and at $R$ one has $Y=X_0$ locally. Consequently $x(s-s')$ has
columns in $R\Lambda$ on the support. Admitted periodicity gives
the same value. If a $P$-neighbor $\Lambda'$ contains $u(X)$,
then the columns of $R s x$ are also in $\Lambda$ because
$sX_0\subset qX$ and $q\Lambda'\subset\Lambda$.
The two lattices agree at $RM$. Admitted locality identifies
$\phi_{\Lambda'}(s x)=c$ there as well.

For each $i$ with $u(X_i^-)\subset\Lambda$, the sum of its
character multipliers over $j$ is either $q^N$ or zero. It is $q^N$
precisely when
\[
                 h(u\eta_i,u(bv))\in\mathcal O
                     \quad\hbox{for all }v\in X ,
\]
where $h$ is the standard definite form, linear in its first argument.
Indeed the character $\chi_H(\beta_{ij};x)$ is
$\exp(2\pi i\operatorname{Tr}_{E/\mathbb Q} h(u\eta_i,u(b\xi_j)))$.
The two terms of the Hermitian matrix supply a CM type and its
conjugate in the trace. Since $QX\subset X_i^-$ and $b\in Q$,
the first vector can have only a simple pole at $Q$ and the second
is integral in $\Lambda$. Self-duality here uses the
$\mathcal O$-valued form. Thus the pairing can have only a simple
pole at $Q$. Replacing $v$ by an element of $PX$ removes that pole
by conjugate linearity. This defines an additive character of $X/PX$.
If it is trivial, scaling $v$ by every element of $\mathcal O$
puts each pairing in the trace dual of $\mathcal O$. At $Q$ this
dual equals $\mathcal O_Q$ since $q$ is prime to the discriminant.
At all other places the pairing was already integral. Conversely
an algebraic integer has integral trace. Summing a character over
the finite additive group proves the assertion.
% END P75H1

% BEGIN P75H2
\paragraph{P75H2. The pointwise Hecke relation, including the nonintegral branch.}
Write $T_Y(\Lambda)= (\mathrm{twist}_R(Y)\phi)_\Lambda(x)$, and
$\Lambda'\sim_P\Lambda$ for the neighbors in Z75D3. With the data
of P75H1 the exact identity is
\[
\begin{split}
 (q-1)\bigg(q^n\sum_i T_{X_i^+}(\Lambda)
       +\sum_{i,j}\chi_H(\beta_{ij};x)T_{X_i^-}(\Lambda)\bigg)
    ={}&(q-1)q^N\sum_{\Lambda'\sim_P\Lambda}T_X(\Lambda')\\
       &+(q^{2N}-q^n)T_X(\Lambda).
\end{split}
\]
It holds even if $c$ of P75H1 vanishes.

If $u(X)\subset\Lambda$, let $e$ be the rank of the induced map
$X/PX\to\Lambda/P\Lambda$. The number of allowed enlargements
is the number of lines in its kernel,
$(q^{N-e}-1)/(q-1)$. Every one of the $A$ hyperplanes contributes
$q^N$ since all pairings are integral. A neighbor containing $u(X)$
corresponds to a hyperplane in $\Lambda/P\Lambda$ containing the
image, hence there are $(q^{n-e}-1)/(q-1)$ of them. These three
counts and P75H1 turn both sides into
$c(q^{n+N-e}-q^n+q^{2N}-q^N)$. This includes the zero map.

Suppose now that $u(X)\not\subset\Lambda$. The enlargements and
the last term are zero. A hyperplane can contribute only if
$X_i^-=u^{-1}(\Lambda)$, since $X_i^-$ has simple quotient in
$X$ and the preimage is proper. In that case the image is integral
away from $Q$ and its pole classes at $Q$ span one nonzero line
$\ell\subset Q^{-1}\Lambda_Q/\Lambda_Q$, generated by
$u\eta_i$. Here $Q^{-1}\Lambda_Q$ denotes multiplication by the
inverse local ideal. Under the perfect pairing with the residue
space at $P$, this line defines a hyperplane $K$ in
$\Lambda_P/P\Lambda_P$. There is just one candidate neighbor:
it shrinks $\Lambda_P$ to the inverse image of $K$, enlarges
$\Lambda_Q$ by $\ell$, and keeps the other completions.
Local duality and the prime-to-level condition give a definite
self-dual neighbor; this is the classification of Z75D3. It contains
the entire image if and only if the $P$-components of $u(v)$ reduce
into $K$ for every $v$. The test for that condition is exactly the
integrality of $h(u\eta_i,u(bv))$ in P75H1. Indeed the only
possible pole is at $Q$, where the second entry uses the conjugate
of the $P$-component, and $b\equiv1\pmod P$.

Conversely any neighbor containing the nonintegral image has a
one-dimensional quotient by its intersection with $\Lambda$, at
$Q$. Pulling that quotient back to $X$ gives exactly a hyperplane
at $Q$ and forces the candidate just described. Thus the two
counts are both zero or both one. Each survivor contributes $q^N c$
to the parenthesis on the left and $c$ to the neighbor sum on the
right. This proves the identity in the second branch as well.
% END P75H2

% BEGIN P75H3
\paragraph{P75H3. The period relation with the original class weights.}
Let $n$ be even, $4N<n$ and $Z$ in the degree-$N$ Hermitian
half space. Write $D_Y(Z)$ for the class average of theta minus
Eisenstein at zero elliptic variables for $\mathrm{twist}_R(Y)\phi$.
On the original Weil transformation laws the pointwise identity gives
\[
 \begin{split}
 (q^N-1)(q^n+q^N)D_X(Z)=(q-1)\bigg(
  q^n\sum_i D_{X_i^+}(Z)
       +\sum_{i,j}D_{X_i^-}(Z+\beta_{ij})\bigg).
 \end{split}
\]
The embedded shifts are understood place by place; each preserves
$\operatorname{Im}Z$. All terms refer to the same admitted family.

To see this, for any family apply the linear functional sending a
test, at a fixed lattice, to the difference of its two series at
$(Z,0,0)$, then sum with the stabilizer weights. The twists have
bounded denominator, period and support as rational Schwartz tests:
choose a common denominator of $Y$ and of its inverse lattice, and
use $u(Y)\subset\Lambda$ and the finite level of $\phi_\Lambda$.
More explicitly the indicator cuts out a full lattice in the matrix
space, and $\phi_\Lambda(sx)$ has a lattice of periods; their
intersection is a lattice of periods for the product, which is
supported on the first lattice. Finite linear combinations and the
rational quadratic characters here remain Schwartz: in integral
coordinates at a common denominator, a sufficiently divisible shift
changes the quadratic exponent by an integer. The defining theta
series is linear by absolute convergence. The Eisenstein sum is
linear by the linear Weil action and stable-range absolute convergence.
The latter estimate is established in the convergence chapter.

For the Hermitian shifts the theta exponential at $Z+\beta$
equals the original exponential times $\chi_H(\beta;x)$.
For Eisenstein use the translation matrix with top right block
$\beta$. Its determinant factor at zero is one. Reindex the bottom
spaces by this rational unitary element. The Weil composition and
the transformation of the test by the inverse translation give
exactly the positive character multiplying the test at $Z$.
This is the same transformation as for theta, so it converts each
twisted character term into $D_{X_i^-}(Z+\beta_{ij})$.

The value at $\Lambda$ of the unaveraged difference for $T_X$ is
invariant under transporting that lattice by a rational unitary
matrix in the defining level group. For the twists this follows
from the equivariance of $\phi$ and the support condition. In
theta the corresponding row change is a reindexing preserving
the quadratic exponent. In Eisenstein the row-unitary action
commutes with the Weil action and preserves evaluation at zero.
Equivalently, reindex the rows in the defining theta transformation:
the two transformed tests have identical theta functions for all
elliptic arguments, whose Fourier coefficients identify the tests.
The zero-argument factors are unchanged. Hence Z75D3 applies to
this particular function, and the class average of its neighbor sum
is $(q^n-1)/(q-1)$ times its original average. Applying the functional
to P75H2 now gives the coefficient
\[
     q^N(q^n-1)+q^{2N}-q^n=(q^N-1)(q^n+q^N)
\]
on $D_X$. This proves the period relation using only the original
character, neighbors and mass.
% END P75H3
